\documentclass{article}
\usepackage{geometry}
\usepackage{graphicx} % Required for inserting images
\usepackage{amsmath, amsthm, amssymb}
\usepackage{parskip}
\usepackage{xr}
\usepackage{amsmath}
\usepackage{amssymb}
\newgeometry{vmargin={15mm}, hmargin={24mm,34mm}}
\theoremstyle{definition} 
\newtheorem{definition}{Definition}

\newtheorem{theorem}{Theorem}[section]
\newtheorem{lemma}[theorem]{Lemma}
\newtheorem{corollary}{Corollary}[theorem]
\newtheorem{proposition}[theorem]{Proposition}
\newtheorem{example}[theorem]{Example}

\newcommand{\N}{\mathbb{N}}
\newcommand{\Z}{\mathbb{Z}}
\newcommand{\R}{\mathbb{R}}
\newcommand{\Q}{\mathbb{Q}}
\newcommand{\C}{\mathbb{C}}
\newcommand{\Nil}{\text{Nil}}
\newcommand{\gauss}{\Z[i]}
\DeclareMathOperator{\Gal}{Gal}
\DeclareMathOperator{\Aut}{Aut}

\title{MATH110C Homework 2}
\date{April 2024}
\author{Boran Erol}

\begin{document}

\maketitle

\section{Exercise 1}

\begin{lemma}
    Let $F$ be a field and $K$ be the field of rational functions
    over $F$. Then, $X^{n} - T$ is irreducible over $K[X]$.
\end{lemma}
\begin{proof}
    Notice that $K[T]$ is the field of fractions of $F[T]$.
    Then, since $F[T]$ is a UFD, we claim that
    Eisenstein's Criterion
    applies with the choice of $T$.

    Notice that $(T)$ is a prime ideal as $F[T]/(T) = F$
    is a domain.
\end{proof}

Then, $\forall n \geq 1: [L : K] \geq n$ since $L$
contains the roots of $X^{n} - T$.

\newpage

\section{Exercise 2}

\begin{lemma}
    Let $\alpha \in \C$ be algebraic over $\Q$
    with degree $n$. Then, $\exists c \in \R: \forall p,q
    \in \Z: q \neq 0: \lvert \alpha - \frac{p}{q} \rvert > \frac{c}{q^{n}}$.
\end{lemma}
\begin{proof}
    Notice that if $\alpha$ has a non-zero imaginary component, we
    can just take $c$ to be the magnitude of the imaginary component
    and trivially satisfy the inequality. Therefore, without loss
    of generality, let $\alpha \in \R$ and $f(x) = a_{n}x^{n} + \cdots + a_{0} \in \Z[x]$ 
    be irreducible. (We can assume it's in $\Z[x]$ by clearing out denominators.)

    Assume that $\lvert \alpha - \frac{p}{q} \rvert < 1$, since otherwise
    we can pick $c \leq 1$ and trivially satisfy the inequality.

    Let $M = \sup_{\lvert x - \alpha \rvert < 1}\{ \lvert f^{\prime}(x) \rvert \}$.

    By the Mean Value Theorem, $\forall p,q
    \in \Z: q \neq 0$, we have that there's some $c$ between $\frac{p}{q}$ and $\alpha$
    such that 

    \[ \lvert f(\frac{p}{q}) - f(\alpha) \rvert \leq \lvert f^{\prime}(c) \rvert \lvert \alpha - \frac{p}{q} \rvert \]
    
    Then,

    \[ \lvert f(\frac{p}{q}) \rvert \leq M \lvert \alpha - \frac{p}{q} \rvert \]

    Since $f(\frac{p}{q}) \neq 0$, $f(\frac{p}{q}) \geq \frac{1}{q^{n}}$. Thus,

    \[ \frac{1}{q^{n}M} \leq \lvert \alpha - \frac{p}{q} \rvert \] 
\end{proof}

\newpage

\section{Exercise 3}



\newpage

\section{Elman pg.311 Exercise 51.5.2}

\begin{definition}
    Let $S$ be a collection of polynomials over $F$. We define the \textbf{splitting field
    of the collection $S$} to be the field where all elements of $S$ factor into linear factors.
\end{definition}

\begin{lemma}
    Let $S$ be a collection of polynomials over $F$. The splitting field of $S$
    exists.
\end{lemma}
\begin{proof}
    Let's first prove existence. Let $K$ be the
    algebraic closure of $F$.
    
    Consider the set of all subfields
    of $K$ that contain all roots of polynomials
    in $S$. This set is non-empty since
    $K$ is in it. Taking the intersection of all
    these sets produces the desired result.
\end{proof}

For uniqueness, I prove an even general result.

\begin{lemma}
    Let $K,L$ be fields and $\sigma: K \xrightarrow{} L$ be an isomorphism.
    Let $S \subseteq K[x]$ and $S^{\prime} \subseteq L[x]$.
    Let $F$ be the splitting field of $S$ over $K$ and $M$
    be the splitting field of $S^{\prime}$ over $L$. Then, we can
    extend $\sigma$ to $\sigma: F \xrightarrow{} M$.
\end{lemma}
\begin{proof}
    Consider $\mathcal{S} = (E,N,\tau)$, where
    $K \subset E \subset F$ and $L \subset N \subset M$
    and $\tau: E \xrightarrow{} N$ extends $\sigma$.

    Notice that the algebraic closures (and the respective
    isomorphism) of these fields are in $\mathcal{S}$.
    Moreover, we can take the union of the fields
    as an upper bound on any chain (identical to how we
    proved uniqueness of algebraic closures in lecture).
    Thus, by Zorn's Lemma, a maximal element of this
    set exists.

    By the exact same argument as how we proved the maximal
    element equals the algebraic closure (namely, using contradiction
    and adding the root of the escaped polynomial), we can
    show that the maximal element of $\mathcal{S}$
    is $(F,M, \tau)$ for some isomorphism $\tau$. 
\end{proof}

\begin{lemma}
    The algebraic closure of $F$ is the splitting field for the collection of irreducible
    polynomials in $F[t]$.
\end{lemma}
\begin{proof}
    Let $K$ be an algebraic closure of $F$. 
    
    We have to show that every irreducible polynomial in $F[t]$
    splits into linear factors over $K$ and that there's at least one irreducible
    polynomial doesn't split over any proper subfield of $K$.

    Let's first show that every irreducible polynomial in $F[t]$ splits over $K$.
    Suppose not. Then, there's a polynomial in $K[t]$ that doesn't split over $K$.
    Thus, $K$ is not algebraically closed, which is a contradiction.

    Let's now show that if every irreducible polynomial in $F[t]$ splits
    over some $E \subseteq K$, $E = K$. Let $E$ be such that 
    every irreducible polynomial in $F[t]$ splits over $E$.
    Then, notice that every polynomial with coefficients in $E$ can be rewritten
    as a polynomial with coefficients in $F$ since $E$ is an algebraic extension over $F$.
    Then, $E$ is algebraically closed. But then $K$ is an algebraic extension of an algebraically
    closed field, and thus $E = K$.

    Let's now prove the converse of the statement, namely, that the splitting field
    for the collection of irreducible polynomials in $F[t]$ is the algebraic closure of $F$.
    
    Let $K$ be the splitting field for the collection of irreducible polynomials in $F[t]$.
    Since every polynomial in $F[t]$ can be written as a product of irreducible polynomials,
    every polynomial in $F[t]$ splits over $K$. Then, it suffices to prove that $K$ is
    an algebraic extension of $F$.

    Let $L$ be the field generated by adding all roots of all polynomials in $F[t]$.
    Since $L$ is algebraic and contains $K$, $K$ is algebraic. We thus conclude the proof.
\end{proof}


\newpage

\section{Elman pg.320 Exercise 52.21.5}

\begin{lemma}
    Let $n \in \Z$. An angle of $2 \pi / n$ radians can be trisected by a straight-edge and compass
    if and only if $3 \nmid n$.
\end{lemma}
\begin{proof}
    Let's first prove the forward direction by using the contrapositive.
    Suppose $n = 3^{a}k$ for some $k \in \Z$ such that $a \geq 1$
    and $\gcd(3,k) = 1$.
    Recall that $[\Q(\frac{2\pi}{n}) : \Q] = \phi(n)$, where $\phi$
    is Euler's Totient Function.

    Then, $[\Q(\frac{2\pi}{3n}) : \Q(\frac{2\pi}{n})] = \phi(3n) / \phi(n)$.

    Now, notice that $\phi(3n) = \phi(3^{a+1}) \phi(k) = 2 \times 3^{a} \phi(k)$
    and $\phi(n) = 2 \times  3^{a-1} \times \phi(k)$, therefore 

    \[ [\Q(\frac{2\pi}{3n}) : \Q(\frac{2\pi}{n})] = 3\].

    We have shown in class that a degree 3 field extension isn't constructible.
  
    Let's now prove the reverse direction. Let $n$ such that $3 \nmid n$.
    Then, $3$ and $n$ are coprime, so by the Bezout Identity, $\exists a,b \in \Z:
    3a + bn = 1$. Then,

    \[ \frac{2\pi}{3n} = \frac{2\pi}{n} a + \frac{2\pi}{3}b\]

    Since $\cos(\frac{2\pi}{3}) = -0.5$, $2\pi/3$ is constructible. Therefore,
    we can trisect $2\pi/n$.    
\end{proof}

\end{document}
